% ============================================================
%  Section 1 — Introduction  (~1.5 pages)
% ============================================================
\section{Introduction}

The quantum software ecosystem has fragmented into a collection of
competing, incompatible SDKs.
Qiskit~\cite{qiskit}, Cirq~\cite{cirq}, and
PennyLane~\cite{bergholm2018pennylane} are the three most widely adopted
frameworks, each offering a distinct programming model and each producing
a different gate decomposition for the same conceptual circuit.
For a practitioner selecting a framework for a NISQ application, this
diversity creates a genuine engineering problem: performance differences
reported in the literature may reflect the choice of simulator as much
as the choice of SDK.

Prior cross-framework comparisons are methodologically flawed in a
consistent way: each framework is executed on its own native simulator.
A difference in measured fidelity could originate from the transpiler,
the simulator, or an interaction between the two.
In the NISQ era---where gate counts and circuit depth directly govern
decoherence and error accumulation~\cite{preskill2018quantum}---this
confound is not a minor caveat; it invalidates the comparison as a guide
to framework selection.

We resolve the confound with QCanvas, a unified compilation platform
that acts as a single intermediate layer between framework source code
and simulation.
The pipeline is: native circuit (Qiskit, Cirq, or PennyLane)
$\to$ QCanvas converter $\to$
OpenQASM 3.0~\cite{cross2022openqasm3}
$\to$ QSim statevector backend.
The backend is fixed for all frameworks; every measured difference is
therefore attributable solely to the framework's transpilation strategy
and gate vocabulary.

This paper makes four contributions:
\begin{enumerate}
  \item A unified benchmarking methodology that eliminates the simulator
        confound via a common OpenQASM 3.0 intermediate representation.
  \item Empirical evaluation of 12 quantum algorithms across 6 metric
        dimensions: compilation time, circuit structure, scaling
        behaviour, simulation performance, statistical equivalence, and
        holistic QPack scores.
  \item Discovery of systematic decomposition incompatibilities in
        PennyLane that produce JSD$=1.000$ divergence at scale, with
        documentation of which algorithms and qubit thresholds trigger
        them.
  \item An open benchmark suite and automated notebook pipeline for
        reproducible framework comparisons.
\end{enumerate}

Section~2 surveys related work.
Section~3 provides background on the three frameworks, OpenQASM 3.0, and
the metrics used.
Section~4 describes the QCanvas pipeline and measurement protocol.
Section~5 presents results across all six metric dimensions.
Section~6 discusses implications for framework selection and identifies
limitations.
Section~7 concludes.
